\documentclass[a4paper,11pt]{article}
\usepackage[utf8]{inputenc}
\usepackage{geometry}
\usepackage{xcolor}
\usepackage{listings}
\usepackage{hyperref}
\usepackage{graphicx}
\usepackage{amsmath}
\usepackage{amssymb}
\usepackage{tcolorbox}
\usepackage{float}
\usepackage{booktabs}

% Page geometry
\geometry{top=2.5cm, bottom=2.5cm, left=2.5cm, right=2.5cm}

% Custom colours
\definecolor{ibmblue}{RGB}{0, 45, 156}
\definecolor{ibmred}{RGB}{250, 75, 76}
\definecolor{purpleaccent}{RGB}{109, 40, 217}
\definecolor{amberaccent}{RGB}{245, 158, 11}
\definecolor{tealaccent}{RGB}{13, 148, 136}
\definecolor{codegreen}{rgb}{0,0.6,0}
\definecolor{codegray}{rgb}{0.5,0.5,0.5}
\definecolor{codepurple}{rgb}{0.58,0,0.82}
\definecolor{backcolour}{rgb}{0.95,0.95,0.92}

% Python code style
\lstdefinestyle{mystyle}{
    backgroundcolor=\color{backcolour},
    commentstyle=\color{codegreen},
    keywordstyle=\color{magenta},
    numberstyle=\tiny\color{codegray},
    stringstyle=\color{codepurple},
    basicstyle=\ttfamily\footnotesize,
    breakatwhitespace=false,
    breaklines=true,
    captionpos=b,
    keepspaces=true,
    numbers=left,
    numbersep=5pt,
    showspaces=false,
    showstringspaces=false,
    showtabs=false,
    tabsize=2
}
\lstset{style=mystyle}

% Custom boxes (titles always in braced form)
\newtcolorbox{studentbox}{colback=blue!5!white,colframe=blue!75!black,title={Student Activity}}
\newtcolorbox{theorybox}{colback=gray!5!white,colframe=gray!50!black,title={Theory Recap}}
\newtcolorbox{warningbox}{colback=yellow!10!white,colframe=orange!80!black,title={Important}}
\newtcolorbox{simulatorbox}{colback=purple!5!white,colframe=purple!75!black,title={Simulator}}
\newtcolorbox{reflectionbox}{colback=teal!5!white,colframe=teal!75!black,title={Reflection Questions}}
\newtcolorbox{connectionbox}{colback=green!5!white,colframe=green!60!black,title={Connections to Previous Labs}}
\newtcolorbox{procedurebox}{colback=red!4!white,colframe=red!70!black,title={Procedure 1 vs Procedure 2}}

% Bra-ket shorthands (guarded: never clash with package-provided macros)
\providecommand{\ket}[1]{|#1\rangle}
\providecommand{\bra}[1]{\langle#1|}
\providecommand{\braket}[2]{\langle#1|#2\rangle}
\providecommand{\ketbra}[2]{|#1\rangle\langle#2|}

\title{\textbf{Laboratory Guide 6: Quantum Key Distribution --- The BB84 Protocol}\\
\large Measurement as a Security Primitive}
\author{Course: Quantum Computing Foundation}
\date{}

\begin{document}

\maketitle

\section{Introduction and Objectives}

In Lab 5 you learned that measuring a qubit in the \emph{wrong} basis gives a random
outcome and \textbf{irreversibly disturbs the state}. That looked like a limitation.
In 1984, Bennett and Brassard turned it into a feature: if an eavesdropper cannot
measure without disturbing, then \textbf{disturbance is evidence of eavesdropping} ---
and two parties can build a shared secret key whose security rests on physics, not on
computational hardness.

This laboratory develops the BB84 protocol completely: the four signal states, the
exact measurement statistics, sifting, the quantum bit error rate (QBER) with and
without an eavesdropper, and the probability of catching her. Every number derived
here is verified live in the lab's interactive simulator and rebuilt with real
quantum circuits in the companion Qiskit notebook.

\vspace{0.3cm}
\noindent\textbf{Prerequisites:}
\begin{itemize}
    \item Lab 5: \textit{Quantum Measurements} --- the Born rule, measurement
          Procedures 1 and 2, and the states $\ket{0}, \ket{1}, \ket{+}, \ket{-}$
          on the Bloch sphere.
    \item \textit{Quantum Communication} (lecture series).
    \item Familiarity with Dirac notation: $\ket{\psi}$, $\bra{\phi}$, $\braket{\phi}{\psi}$.
\end{itemize}

\begin{connectionbox}
\textbf{The Lab 5 bridge.} Everything in BB84 is applied Lab 5:
\begin{itemize}
    \item The four signal states are exactly the states $\ket{0}, \ket{1}, \ket{+},
          \ket{-}$ you placed on the Bloch sphere --- the poles and two equator points.
    \item The measurement statistics are the Born rule you derived, evaluated at those
          four states in the $Z$- and $X$-bases (Part B).
    \item The distinction between \textbf{Procedure 1} (direct basis projection) and
          \textbf{Procedure 2} (rotate with the canonical rotation $U_X = H$, then
          $Z$-measure) --- which may have seemed academic --- becomes \emph{operationally
          decisive} when the eavesdropper must resend a quantum state (Part D).
\end{itemize}
\end{connectionbox}

\noindent\textbf{Mode:} Interactive Simulator + Qiskit \quad\textbf{Credits:} 1 ECTS
\quad\textbf{Estimated time:} $\approx$ 10 h

\vspace{0.3cm}
\noindent\textbf{Learning Objectives:}
\begin{enumerate}
    \item \textbf{Explain the logic of BB84 and its security foundations} --- random
          bases, mutually unbiased states, sifting, and why
          $|\braket{z_j}{x_l}|^2 = 1/2$ makes a wrong-basis measurement erase the
          encoded bit.
    \item \textbf{Perform and analyze QKD experiments} --- run the protocol in the
          simulator and in Qiskit, compute the sifted fraction and the QBER, and
          compare against the exact theoretical values.
    \item \textbf{Understand the measurement--disturbance--security relationship} ---
          derive the 25\% QBER of intercept-resend, see why the eavesdropper cannot
          avoid it, and quantify detection: $P_{\text{det}}(k) = 1 - (3/4)^k$.
    \item \textbf{Discuss the relevance for cryptographic systems} --- key rates,
          error thresholds, and what QKD does (and does not) replace in modern
          cryptography.
\end{enumerate}

\newpage

\section{Part A: The Four BB84 States}

Alice encodes each random bit $a \in \{0,1\}$ in one of two randomly chosen bases:
the \textbf{$Z$-basis} (rectilinear) or the \textbf{$X$-basis} (diagonal). The
encoding rule: bit 0 always maps to the $+1$ eigenstate of the basis observable,
bit 1 to the $-1$ eigenstate.

\begin{equation}
    \ket{\psi(0,Z)} = \ket{0}, \qquad
    \ket{\psi(1,Z)} = \ket{1}, \qquad
    \ket{\psi(0,X)} = \ket{+}, \qquad
    \ket{\psi(1,X)} = \ket{-},
    \label{eq:encoding}
\end{equation}

\noindent with $\ket{\pm} = (\ket{0} \pm \ket{1})/\sqrt{2}$. On the Lab 5 Bloch
sphere these are the poles ($\theta = 0, \pi$) and two equator points
($\theta = \pi/2$, $\varphi = 0, \pi$). In the photon-polarization picture they
correspond to vertical, horizontal, and the two diagonal polarizations.

Both bases are orthonormal, and they are \textbf{mutually unbiased}. Writing
$\ket{z_0} = \ket{0}$, $\ket{z_1} = \ket{1}$, $\ket{x_0} = \ket{+}$,
$\ket{x_1} = \ket{-}$:

\begin{equation}
    |\braket{z_j}{x_l}|^2 = \frac{1}{2}
    \qquad \text{for all } j, l \in \{0, 1\}.
    \label{eq:mub}
\end{equation}

\noindent Direct computation: $\braket{0}{\pm} = 1/\sqrt{2}$ and
$\braket{1}{\pm} = \pm 1/\sqrt{2}$, so every cross-basis overlap has squared
modulus $1/2$.

\begin{theorybox}
\textbf{Why this is the security foundation.} A measurement in the wrong basis
yields a uniformly random outcome --- it extracts \emph{zero} information about the
encoded bit and (Lab 5!) destroys the original state in the process. Everything
that follows is this one fact, exploited systematically.
\end{theorybox}

\section{Part B: Measurement Statistics --- the Born Table}

Applying the Born rule from Lab 5,
\begin{equation}
    P(m \,|\, \psi; B) = |\braket{m}{\psi}|^2,
    \label{eq:born}
\end{equation}
to the four BB84 states in both bases gives the complete statistics of the protocol:

\begin{table}[H]
\centering
\begin{tabular}{lcccc}
\toprule
 & \multicolumn{2}{c}{\textbf{$Z$-measurement}} & \multicolumn{2}{c}{\textbf{$X$-measurement}} \\
\textbf{State sent} & $P(0)$ & $P(1)$ & $P(+)$ & $P(-)$ \\
\midrule
$\ket{0}$  & $1$   & $0$   & $1/2$ & $1/2$ \\
$\ket{1}$  & $0$   & $1$   & $1/2$ & $1/2$ \\
$\ket{+}$  & $1/2$ & $1/2$ & $1$   & $0$   \\
$\ket{-}$  & $1/2$ & $1/2$ & $0$   & $1$   \\
\bottomrule
\end{tabular}
\caption{The BB84 Born table. Matched basis $\Rightarrow$ deterministic recovery of
the bit; mismatched basis $\Rightarrow$ coin-flip outcome, bit erased
(equation~\eqref{eq:mub}).}
\label{tab:born}
\end{table}

\begin{procedurebox}
\textbf{Read this carefully --- it returns in Part D.} Table~\ref{tab:born} lists
\emph{probabilities}, and these are identical whether the $X$-basis measurement is
done by \textbf{Procedure 1} (direct projection onto $\{\ket{+}, \ket{-}\}$) or by
\textbf{Procedure 2} (apply the canonical rotation $U_X = H$, then $Z$-measure ---
what real hardware does). Lab 5 established this equivalence. Bob keeps only a
classical bit, so \emph{for Bob} the choice of procedure is invisible: with
Procedure 2 the $Z$-outcome $c \in \{0,1\}$ maps to key bit $b = c$ because
$H\ket{+} = \ket{0}$ and $H\ket{-} = \ket{1}$.

\medskip
But the two procedures leave \textbf{different post-measurement states}
(Procedure 1: a basis eigenstate; Procedure 2: $\ket{0}$ or $\ket{1}$). Bob never
cares. \textbf{The eavesdropper will} --- she must resend a quantum state, and a
post-measurement state is exactly what she has in hand. Hold that thought for
Part D.
\end{procedurebox}

\section{Part C: The Protocol, Sifting, and the Error Check}

For each round $i = 1, \dots, n$:

\begin{enumerate}
    \item \textbf{Alice} draws a uniform random bit $a_i$ and basis
          $\alpha_i \in \{Z, X\}$, prepares $\ket{\psi(a_i, \alpha_i)}$, and sends
          it through the quantum channel.
    \item \textbf{Bob} draws his own uniform random basis $\beta_i \in \{Z, X\}$ ---
          independently, since he cannot know $\alpha_i$ --- and measures, recording
          bit $b_i$. (On hardware: Procedure 2 with $U_X = H$.)
    \item \textbf{Basis reconciliation:} over an authenticated public channel, Alice
          and Bob announce their \emph{bases} --- never the bits. They keep the
          \textbf{sifted set} $S = \{i : \alpha_i = \beta_i\}$ and discard everything
          else.
    \item \textbf{Error check:} they sacrifice a uniformly random subset of $k$
          sifted positions, publicly compare $a_i$ vs $b_i$ there, and estimate the
          \textbf{quantum bit error rate (QBER)}. Too many errors $\rightarrow$ abort.
    \item \textbf{Key:} the remaining sifted, unchecked bits form the raw shared key.
\end{enumerate}

\begin{theorybox}
\textbf{Sifted fraction.} $\alpha_i$ and $\beta_i$ are independent and uniform on
two values, so
\begin{equation}
    P(\alpha_i = \beta_i) = P(Z,Z) + P(X,X) = \tfrac{1}{4} + \tfrac{1}{4}
    = \tfrac{1}{2}.
    \label{eq:sift}
\end{equation}
On average \textbf{half the rounds survive sifting}: the sifted length is
binomially distributed with expectation $n/2$.
\end{theorybox}

\begin{theorybox}
\textbf{QBER without eavesdropping.} On an ideal channel, every sifted round is a
matched-basis measurement --- the deterministic cells of Table~\ref{tab:born}. Hence
\begin{equation}
    Q_{\text{no Eve}} = P(b_i \neq a_i \mid i \in S) = 0 \quad \text{exactly.}
    \label{eq:qzero}
\end{equation}
Not ``small'' --- \textbf{zero}. On an ideal channel, \emph{any} error in the sifted
key is evidence that something interfered with the qubits in transit.
\end{theorybox}

\begin{warningbox}
\textbf{Why sifting matters.} Without step 3, mismatched-basis rounds (coin flips)
would contribute errors even with no eavesdropper at all: the raw, unsifted
disagreement rate is $\tfrac{1}{2} \cdot \tfrac{1}{2} = 25\%$ for a perfectly
secure channel. Sifting removes exactly this noise floor, which is what makes the
QBER a meaningful security signal. Verify both numbers in the simulator.
\end{warningbox}

\section{Part D: The Intercept-Resend Attack --- QBER = 25\%}

Eve's simplest attack: intercept each qubit, measure it in a random basis
$\gamma_i \in \{Z, X\}$, and resend the post-measurement state to Bob. She cannot
do better than guess --- the basis announcements happen only \emph{after} Bob has
measured. Condition on a sifted round ($\beta = \alpha$) and split on Eve's guess:

\begin{itemize}
    \item \textbf{Case 1 --- Eve guesses right} ($\gamma = \alpha$, probability
          $\tfrac{1}{2}$). Matched-basis measurement is deterministic: Eve reads
          $e = a$, resends the very same state, Bob recovers $b = a$.
          Error probability: $0$.
    \item \textbf{Case 2 --- Eve guesses wrong} ($\gamma \neq \alpha$, probability
          $\tfrac{1}{2}$). Her outcome is a coin flip, and her resent state is a
          $\gamma$-basis eigenstate. Bob measures it in $\alpha \neq \gamma$ ---
          mutually unbiased again --- so his outcome is uniform \emph{regardless of
          what Alice sent}. Error probability: $\tfrac{1}{2}$.
\end{itemize}

\begin{equation}
    Q \;=\; \underbrace{\tfrac{1}{2} \cdot 0}_{\gamma = \alpha}
       \;+\; \underbrace{\tfrac{1}{2} \cdot \tfrac{1}{2}}_{\gamma \neq \alpha}
       \;=\; \tfrac{1}{4} \;=\; 25\%.
    \label{eq:qber25}
\end{equation}

\noindent Two refinements, both testable directly in the simulator:

\begin{itemize}
    \item \textbf{Partial interception.} If Eve intercepts each round independently
          with probability $p$ and lets the rest pass, errors arise only on
          intercepted rounds:
          \begin{equation}
              Q(p) = \frac{p}{4}, \qquad Q(0) = 0, \qquad Q(1) = \tfrac{1}{4},
              \label{eq:partial}
          \end{equation}
          linear in $p$. Less interception means less information for Eve \emph{and}
          less visibility --- but never zero of one with the other.
    \item \textbf{Fixed-basis Eve.} If Eve always measures in $Z$ instead of
          guessing randomly, the arithmetic changes but the answer does not: her
          matching rounds contribute $0$, the others $\tfrac{1}{2}$ --- still
          $Q = \tfrac{1}{4}$. Randomizing her basis neither helps nor hurts her
          visibility.
\end{itemize}

\begin{procedurebox}
\textbf{Where the procedure distinction becomes operational.} ``Resend the
post-measurement state'' hides a subtlety. If Eve measures by \textbf{Procedure 1},
her post-measurement state \emph{is} the $\gamma$-basis eigenstate
$\ket{e_\gamma}$ --- resending it gives the 25\% above. If she measures by
\textbf{Procedure 2} (apply $U_\gamma$, $Z$-measure), her post-measurement state is
the computational state $\ket{e}$, and to run the same attack she must
\textbf{re-encode} it by applying $U_\gamma^\dagger$ before resending. Done
correctly, the two implementations are indistinguishable --- same states to Bob,
same $Q = 25\%$.

\medskip
\textbf{But suppose she forgets the re-encoding} and naively resends $\ket{e}$
as-is. Work it through with Table~\ref{tab:born}:
\begin{itemize}
    \item Alice's $Z$-rounds: Eve's $\gamma = Z$ causes no error; $\gamma = X$
          scrambles $e$ to a coin flip that Bob then reads deterministically
          $\rightarrow$ error $\tfrac{1}{2}$. Contribution:
          $\tfrac{1}{2} \cdot \tfrac{1}{2} = \tfrac{1}{4}$.
    \item Alice's $X$-rounds: Bob $X$-measures a $Z$-eigenstate --- coin flip
          \emph{whatever} Eve's basis was $\rightarrow$ error $\tfrac{1}{2}$.
          Contribution: $\tfrac{1}{2}$.
\end{itemize}
\begin{equation}
    Q_{\text{naive}} = \tfrac{1}{2} \cdot \tfrac{1}{4}
                     + \tfrac{1}{2} \cdot \tfrac{1}{2}
                     = \tfrac{3}{8} = 37.5\% \;\neq\; 25\%.
    \label{eq:naive}
\end{equation}
Every outcome \emph{probability} in her measurement was identical between the two
procedures --- yet the attack statistics changed, because the
\emph{post-measurement states} differ and the resent state is a post-measurement
state. This is the Lab 5 distinction made operational. The simulator's ``naive
resend'' mode lets you watch the QBER converge to 37.5\% instead of 25\%.
\end{procedurebox}

\section{Part E: Catching Eve --- Detection Probability}

Under full interception, each of the $k$ publicly compared check bits is wrong
independently with probability $\tfrac{1}{4}$. Eve escapes only if \emph{all} $k$
agree:

\begin{equation}
    P(\text{undetected}) = \left(\tfrac{3}{4}\right)^k
    \qquad\Longrightarrow\qquad
    P_{\text{det}}(k) = 1 - \left(\tfrac{3}{4}\right)^k.
    \label{eq:detect}
\end{equation}

\begin{table}[H]
\centering
\begin{tabular}{lcccc}
\toprule
check bits $k$      & 1      & 5      & 10     & 20     \\
\midrule
$P_{\text{det}}(k)$ & $0.25$ & $\approx 0.7627$ & $\approx 0.9437$ & $\approx 0.9968$ \\
\bottomrule
\end{tabular}
\caption{Detection probability against full intercept-resend,
equation~\eqref{eq:detect}.}
\label{tab:detect}
\end{table}

\noindent Detection is \textbf{exponentially certain} in the number of check bits.
For partial interception the same argument gives
\begin{equation}
    P_{\text{det}}(k; p) = 1 - \left(1 - \frac{p}{4}\right)^k
    \label{eq:detectp}
\end{equation}
(e.g.\ $p = 0.5$, $k = 16$ gives $P_{\text{det}} \approx 0.8819$). Note the
$k = 0$ limit: no check bits, no detection --- the public comparison is not an
optional optimization, it \emph{is} the security test.

\begin{warningbox}
\textbf{Honest security context.} What you derived here covers one specific attack.
The full security statement for BB84 (Shor \& Preskill, 2000) says the protocol
with error correction and privacy amplification produces a secure key whenever the
QBER stays below $\approx 11\%$ --- against \emph{any} attack physics allows, not
just intercept-resend. Intercept-resend's 25\% sits far above that threshold, which
is why this attack is always catastrophic for Eve. The proof machinery is beyond
this lab; the QBER-threshold \emph{logic} --- measure the error rate, abort above
threshold --- is exactly what you practice here.
\end{warningbox}

\newpage

\section{Part F: Guided Activities}

\begin{simulatorbox}
The lab's interactive page (\texttt{lab6\_bb84.html}) contains a BB84 simulator
built for this lab. \textbf{Step-by-Step} mode shows one photon at a time --- watch
the basis choices, Eve's interference, and the sifted key forming.
\textbf{Batch Statistics} mode runs thousands of rounds and plots the measured
statistics against every formula derived above. All runs are seeded: the same seed
always reproduces the same run. The companion Qiskit notebook rebuilds the same
protocol with real quantum circuits:\\[4pt]
\url{https://colab.research.google.com/github/LMrilo/quantum-labs-2026/blob/main/lab6/lab6_bb84.ipynb}
\end{simulatorbox}

\subsection{Activity 1: Sifting Statistics (Batch mode)}

\begin{studentbox}
With Eve \textbf{off}, run $N = 100$, then $N = 1000$, then $N = 5000$ (same seed).
Record the sifted fraction each time.
\begin{enumerate}
    \item How does the deviation from $1/2$ shrink as $N$ grows? Compare with the
          statistical scale $\sqrt{1/(4N)}$.
    \item Confirm the QBER is \emph{exactly} 0 in every run --- not just small. Why
          is this exact even though the simulation is random?
          (Equation~\eqref{eq:qzero}.)
\end{enumerate}
\end{studentbox}

\subsection{Activity 2: Measuring Eve (Step-by-Step + Batch modes)}

\begin{studentbox}
Turn Eve \textbf{on} (intercept-resend, $p = 1$).
\begin{enumerate}
    \item In Step-by-Step mode, find a round where Eve guessed the wrong basis but
          \emph{no error resulted}. Why is that possible? (Hint: Case 2 gives an
          error with probability $\tfrac{1}{2}$, not $1$.)
    \item In Batch mode with $N = 5000$, verify the QBER lands near 25\%. How large
          is the deviation? Change the seed a few times --- does it stay within a
          band of roughly $\pm 2\sqrt{0.25 \cdot 0.75/(N/2)}$?
\end{enumerate}
\end{studentbox}

\subsection{Activity 3: The Naive Resend --- Post-States Matter}

\begin{studentbox}
Switch the Eve strategy to \textbf{naive resend} (she measures via Procedure 2 and
forgets to re-encode with $U_\gamma^\dagger$). Run $N = 5000$.
\begin{enumerate}
    \item Verify the QBER converges near 37.5\%, not 25\%
          (equation~\eqref{eq:naive}).
    \item Every measurement \emph{probability} in Eve's apparatus is identical in
          the two strategies. What, physically, is different about what Bob
          receives?
    \item Alice's $Z$-rounds and $X$-rounds contribute differently ($\tfrac{1}{4}$
          vs $\tfrac{1}{2}$). Which rounds in the Step-by-Step log show it?
\end{enumerate}
\end{studentbox}

\subsection{Activity 4: Detection Budget (Detection panel)}

\begin{studentbox}
Suppose Alice and Bob want to catch a full intercept-resend Eve with at least 99\%
probability.
\begin{enumerate}
    \item From $P_{\text{det}}(k) = 1 - (3/4)^k$, compute the smallest sufficient
          $k$. Verify with the detection panel.
    \item Repeat for a lazier Eve with $p = 0.25$
          (equation~\eqref{eq:detectp}). How many check bits do you need now? What
          does this say about the cost of detecting weak eavesdropping?
\end{enumerate}
\end{studentbox}

\subsection{Activity 5: The Protocol in Qiskit}

The companion notebook implements the four BB84 states and the basis-selective
measurement as two short routines, then scales them to the full protocol with
sifting, an intercept-resend Eve (measure-and-re-prepare), QBER statistics against
the 25\% theory line, and the detection probability against $1 - (3/4)^k$. Its two
core routines:

\begin{lstlisting}[language=Python]
def prepare_bb84_state(bit: int, basis: str) -> QuantumCircuit:
    """Return a 1-qubit circuit preparing the BB84 state |psi(bit, basis)>."""
    qc = QuantumCircuit(1)
    if bit == 1:
        qc.x(0)          # |0> -> |1>
    if basis == "X":
        qc.h(0)          # |0> -> |+>,  |1> -> |->
    return qc

def measure_in_basis(qc: QuantumCircuit, basis: str) -> QuantumCircuit:
    """Append a basis-B measurement (Procedure 2) to a 1-qubit circuit, in place."""
    if basis == "X":
        qc.h(0)          # canonical rotation U_X = H
    qc.measure_all()
    return qc
\end{lstlisting}

\noindent The $X$-basis measurement is Procedure 2 with the canonical rotation
$U_X = H$ --- hardware has no other option. The notebook closes with an optional
run on real IBM hardware, where you will meet a nonzero QBER \emph{without} any
Eve, and confront the question real QKD systems face: how do you tell noise from
eavesdropping?

\begin{warningbox}
\textbf{Qiskit bit-ordering note.} Qiskit orders multi-qubit bitstrings
little-endian (rightmost bit = qubit $q_0$), while this lab series writes states
big-endian (qubit 0 leftmost). In this lab every quantum register holds a
\emph{single} qubit, so the two conventions coincide and no conversion is needed
--- stated here for consistency with the multi-qubit labs, where the distinction
is essential.
\end{warningbox}

\subsection{Activity 6: Quick Check}

\begin{reflectionbox}
\textbf{Question 1:} Alice sends $\ket{+}$ and Bob measures in the $Z$-basis. What
is the probability he reads 0?
\begin{itemize}
    \item[(a)] $1$ \qquad (b) $1/2$ \qquad (c) $0$
\end{itemize}

\textbf{Question 2:} Why does the sifted key contain \emph{no} errors on an ideal
channel without Eve?
\begin{itemize}
    \item[(a)] Because errors are corrected during basis reconciliation.
    \item[(b)] Because matched-basis measurements reproduce the encoded bit with
               probability 1.
    \item[(c)] Because Bob's random basis averages the errors away.
\end{itemize}

\textbf{Question 3:} Eve measures a qubit in the $X$-basis via Procedure 2 ($H$,
then $Z$-measure) and gets outcome 0. Which state does she hold afterwards?
\begin{itemize}
    \item[(a)] $\ket{+}$ \qquad (b) $\ket{0}$ \qquad (c) Whatever Alice sent.
\end{itemize}
\end{reflectionbox}

\section{Resources}

\begin{itemize}
    \item Bennett, C. H. \& Brassard, G., ``Quantum cryptography: Public key
          distribution and coin tossing,'' \textit{Proc.\ IEEE Int.\ Conf.\ on
          Computers, Systems and Signal Processing}, Bangalore, 1984, pp.~175--179.
          Reprinted: \textit{Theor.\ Comput.\ Sci.}\ \textbf{560}, 7--11 (2014).
          --- The BB84 protocol.
    \item Nielsen, M. A. \& Chuang, I. L., \textit{Quantum Computation and Quantum
          Information}, Cambridge University Press, 10th anniversary edition, 2010,
          \S 12.6 (quantum cryptography: BB84, intercept-resend, security sketch).
    \item Gisin, N., Ribordy, G., Tittel, W. \& Zbinden, H., ``Quantum
          cryptography,'' \textit{Rev.\ Mod.\ Phys.}\ \textbf{74}, 145 (2002). ---
          Survey; QBER terminology and practical attack context.
    \item Shor, P. W. \& Preskill, J., ``Simple proof of security of the BB84
          quantum key distribution protocol,'' \textit{Phys.\ Rev.\ Lett.}\
          \textbf{85}, 441 (2000). --- The $\approx 11\%$ threshold cited in Part E.
    \item Lab 5: \textit{Quantum Measurements} (this series) --- Born rule,
          Procedures 1 and 2, the states $\ket{0}, \ket{1}, \ket{+}, \ket{-}$.
    \item Companion Jupyter notebook:
          \texttt{quantum-labs-2026/lab6/lab6\_bb84.ipynb} (Colab link in Part F).
\end{itemize}

\end{document}
